\chapter{Introdução}
\label{cap:introducao}

Organizações dependem de processos documentados para operar de forma previsível, treinar pessoas e cumprir exigências regulatórias. A notação BPMN 2.0, padronizada pela \citeonline{omg2011bpmn}, tornou-se a linguagem de referência para essa documentação, com uso consolidado tanto na indústria quanto na literatura de gestão de processos \cite{weske2019business, dumas2018fundamentals}. Na prática, porém, a maior parte do conhecimento sobre processos permanece registrada em linguagem natural: manuais internos, descrições de procedimentos e documentos de política. A conversão desse material em modelos formais continua sendo uma tarefa manual, lenta e dependente de especialistas \cite{friedrich2011process, vanderaa2018challenges}.

O avanço dos modelos de linguagem de grande porte \cite{vaswani2017attention, brown2020language} tornou plausível automatizar essa conversão. A dificuldade está na forma de saída. A serialização XML da BPMN foi concebida para intercâmbio entre ferramentas, não para ser escrita diretamente por pessoas ou por modelos. Ao solicitar que um modelo produza XML BPMN completo, transfere-se a ele uma responsabilidade pouco adequada ao seu funcionamento: reproduzir com precisão um grande volume de marcação técnica, ainda que boa parte dessa marcação possa ser derivada de maneira determinística. O custo aparece em duas frentes. A primeira é econômica e não depende de verificação: cada \textit{token} de marcação repetitiva consome contexto, tempo e dinheiro, e a quantidade emitida é diretamente mensurável. A segunda é de confiabilidade, e tem natureza distinta --- é uma \textbf{expectativa a ser testada}, não um dado. O risco existe em princípio, pois uma única inconsistência estrutural invalida o documento inteiro; se ele de fato se materializa, e em que medida, é questão empírica. Adianta-se que a resposta obtida neste trabalho contraria a expectativa inicial, conforme discutido adiante.

A redução do número de \textit{tokens} em sistemas baseados em modelos de linguagem tem sido atacada principalmente do lado da entrada. Na pesquisa, o LLMLingua \cite{jiang2023llmlingua} comprime o contexto antes que ele alcance o modelo, com ganhos expressivos e pouca perda de desempenho. Na indústria, ferramentas como o Headroom \cite{headroom2026} aplicam a mesma ideia a saídas de ferramentas, registros e trechos recuperados, relatando reduções entre 60\% e 95\% em conteúdo estruturado. Essas abordagens confirmam que a economia de \textit{tokens} é uma preocupação concreta e que boa parte dela vem da remoção de repetição estrutural. Nenhuma delas, entretanto, atua sobre o lado da geração: fazer com que o próprio modelo emita nativamente uma representação compacta, cuja expansão determinística produza um artefato válido segundo um esquema formal.

Duas linhas de trabalho se aproximam desse objetivo por caminhos distintos. A decodificação restrita, discutida por \citeonline{poesia2022synchromesh}, garante conformidade sintática durante a geração, mas não reduz o volume que o modelo precisa emitir. Já o ProMoAI, de \citeonline{kourani2024promoai}, gera modelos de processo por meio de uma representação intermediária expandida deterministicamente em BPMN, demonstrando que a estratégia é viável. O trabalho de \citeonline{kourani2024promoai} apoia-se, contudo, em um modelo de fronteira operado por \textit{prompt} e não trata a representação intermediária como um problema de compressão com taxa medida. É exatamente nesse espaço que esta monografia se posiciona.

A escolha do modelo também é parte do problema. \citeonline{belcak2025small} argumentam que modelos pequenos e especializados, e não modelos de fronteira genéricos, tendem a ser a base adequada para tarefas agênticas repetitivas, por razões de custo, latência e previsibilidade. Uma representação intermediária suficientemente compacta e regular torna essa especialização viável: reduz o que o modelo precisa aprender a emitir e permite que um modelo de sete bilhões de parâmetros, como o Qwen2.5-Coder \cite{hui2024qwen25coder}, seja ajustado para a tarefa.

Este trabalho propõe, portanto, um protocolo de compressão semântica para geração de BPMN. Um modelo de linguagem produz uma descrição concisa do processo em uma linguagem de domínio específico projetada para esse fim, e um transpilador baseado em gramática EBNF expande essa descrição em XML BPMN 2.0 validado contra o esquema oficial. A separação de responsabilidades traz um efeito colateral útil para a depuração: quando ocorre uma falha, ela fica localizada na geração da DSL, no transpilador ou na validação, em vez de diluída em um documento XML extenso. Medido sobre 1.021 processos com o tokenizador do Qwen2.5-Coder, o artefato proposto é \textbf{seis vezes} mais compacto que o XML lógico correspondente, o que equivale a uma redução de 83,4\% no número de \textit{tokens} que o modelo precisa emitir; sobre as saídas efetivamente geradas pelos braços experimentais, a razão sobe para cerca de sete.

Os resultados obtidos são antecipados aqui, por serem em parte contrários à expectativa que motivou o trabalho. A economia confirma-se sem ressalva. A confiabilidade, porém, depende de \textbf{como} a representação intermediária chega ao modelo. Descrita por \textit{prompt}, ela piora o desempenho: modelos de fronteira instruídos pela gramática produziram saída válida em 64,8\% e 56,6\% das gerações, contra 98,1\% e 94,3\% da geração direta de XML pelos mesmos modelos. Internalizada nos pesos por ajuste fino, ela melhora substancialmente: um modelo de sete bilhões de parâmetros atingiu 86,8\% de validade, vinte e dois pontos percentuais acima do modelo de fronteira instruído pela mesma gramática, embora ainda abaixo da geração direta. A explicação é que a serialização XML da BPMN, embora verbosa, é amplamente representada no pré-treinamento dos modelos, ao passo que a linguagem aqui proposta lhes é inteiramente nova --- verbosidade não constitui dificuldade para quem já viu o formato, novidade sim. Segue-se que o ajuste fino não é etapa acessória da abordagem, e sim condição de sua viabilidade.

Registra-se ainda um achado de ordem metodológica, obtido durante a avaliação e independente da linguagem proposta: mediu-se a concordância entre modelos de referência produzidos por especialistas distintos para o mesmo processo, e ela é baixa. Métricas estruturais que identificam nós por rótulo saturam, portanto, muito abaixo do máximo teórico, e todos os métodos avaliados operam nessa vizinhança. A consequência para a área é que valores absolutos dessas métricas não devem ser lidos como percentual de acerto.

\section{Motivação}
\label{sec:motivacao}

A modelagem de processos permanece cara justamente onde é mais necessária. Organizações que já documentam seus procedimentos em texto raramente dispõem de tempo especializado para convertê-los em modelos formais, e o resultado é uma base documental que não pode ser analisada, simulada ou automatizada. Reduzir o custo dessa conversão amplia o acesso à modelagem formal para além das organizações que podem manter analistas dedicados.

Há ainda uma motivação técnica, e ela independe do desfecho particular observado na BPMN. O padrão de separar a geração de uma representação compacta da expansão determinística em um artefato validável é aplicável a qualquer domínio em que a saída precise obedecer a um esquema formal e a serialização oficial seja verbosa. O que este trabalho tem a oferecer a esses domínios não é apenas a confirmação de um ganho, mas a identificação da \textbf{condição} sob a qual ele se realiza --- questão que só uma avaliação com protocolo previamente fixado poderia responder em qualquer direção. A BPMN oferece um caso de teste favorável para essa investigação porque possui esquema oficial, conjuntos de dados públicos com modelos de referência e métricas estruturais bem definidas, o que permite verificação objetiva em vez de avaliação subjetiva.

Por fim, a economia de \textit{tokens} deixou de ser detalhe de implementação. Em fluxos que geram modelos em escala, o custo de emitir marcação redundante é recorrente, e a alternativa de depender continuamente de modelos de fronteira concentra esse custo em um fornecedor externo. Um modelo pequeno ajustado para emitir uma representação comprimida é, nesse cenário, uma resposta de engenharia e não apenas um exercício acadêmico.

\section{Objetivos}
\label{sec:objetivos}

\subsection{Objetivo Geral}
\label{sec:objetivo-geral}

Investigar se a geração de uma representação intermediária comprimida, seguida de expansão determinística, produz modelos BPMN 2.0 válidos de forma mais confiável e mais econômica em \textit{tokens} do que a geração direta de XML por modelos de linguagem.

\subsection{Objetivos Específicos}
\label{sec:objetivos-especificos}

	\begin{alineas}
		\item projetar uma linguagem de domínio específico para descrição de processos BPMN, com gramática EBNF formalmente especificada;
		\item implementar um transpilador determinístico que expanda essa linguagem em XML BPMN 2.0 validado contra o esquema oficial;
		\item construir um conjunto de dados de pares formados por descrição textual e representação na linguagem proposta, a partir de fontes públicas de descrições de processos;
		\item realizar o ajuste fino supervisionado de um modelo de linguagem de pequeno porte para emitir a representação comprimida;
		\item quantificar a compressão obtida e a fidelidade estrutural dos modelos gerados, comparando a abordagem proposta com a geração direta de XML sob protocolo de avaliação definido previamente à execução dos experimentos.
	\end{alineas}

\section{Organização do Trabalho}
\label{sec:organizacao}

O Capítulo~\ref{cap:fundamentacao-teorica} reúne os conceitos necessários à proposta: modelagem de processos e a notação BPMN 2.0, funcionamento e limitações dos modelos de linguagem na geração estruturada, linguagens de domínio específico como camada de mediação, e os métodos de adaptação de modelos empregados na etapa experimental.

O Capítulo~\ref{cap:trabalhos-relacionados} posiciona a proposta frente à literatura, distinguindo-a tanto das abordagens de compressão de contexto, que atuam sobre a entrada, quanto dos trabalhos que geram modelos de processo por representação intermediária sem tratá-la como problema de compressão com taxa medida.

O Capítulo~\ref{chap:metodologia} descreve a linguagem proposta e sua gramática, o transpilador determinístico e o gerador de leiaute, o \textit{pipeline} de construção do conjunto de dados, e --- de forma destacada --- o protocolo de avaliação, congelado antes da execução de qualquer experimento e reproduzido ali na íntegra, incluindo hipóteses, métricas, braços experimentais e plano de análise estatística.

O Capítulo~\ref{chap:resultados} reporta a verificação do componente determinístico sobre o corpus completo e, em seguida, a avaliação comparativa dos sete braços experimentais contra modelos de referência produzidos por especialistas. Reporta também as emendas ao protocolo, as verificações de validade aplicadas aos resultados adversos e as limitações identificadas na própria métrica de fidelidade.

O Capítulo~\ref{chap:conclusoes-e-trabalhos-futuros} retoma o objetivo geral à luz das evidências, sistematiza as contribuições e as limitações, e enumera as extensões que a investigação deixou em aberto.
